\documentclass{tietreport}

% ==================================================
% PACKAGES
% ==================================================

\usepackage{graphicx}
\usepackage{float}
\usepackage{array}
\usepackage{booktabs}
\usepackage{titlesec}

% Fix fancyhdr warning
\setlength{\headheight}{15pt}

% Prevent unwanted blank pages between chapters
\let\cleardoublepage\clearpage

% Avoid excessive vertical stretching
\raggedbottom

% Slightly reduce heading spacing
\titlespacing*{\chapter}
{0pt}
{-20pt}
{20pt}

\titlespacing*{\section}
{0pt}
{14pt}
{6pt}

\titlespacing*{\subsection}
{0pt}
{10pt}
{4pt}

% Reduce figure spacing
\setlength{\textfloatsep}{10pt}
\setlength{\floatsep}{8pt}
\setlength{\intextsep}{8pt}


% ==================================================
% DIAGRAM COMMAND
% ==================================================

% Large diagrams, centred and allowed to extend
% slightly into the page margins.
\newcommand{\diagramfigure}[3]{%
\begin{figure}[H]
    \centering
    \IfFileExists{#1}{%
        \makebox[\textwidth][c]{%
            \includegraphics[
                width=1.10\textwidth,
                height=0.78\textheight,
                keepaspectratio
            ]{#1}
        }
    }{%
        \fbox{%
            \parbox[c][5cm][c]{0.9\textwidth}{%
                \centering
                Diagram not yet uploaded:\\[2mm]
                \texttt{#1}
            }
        }
    }
    \caption{#2}
    \label{#3}
\end{figure}
}


% ==================================================
% PROJECT INFORMATION
% ==================================================

\TitlePageHeader{%
    UCS503P - Software Engineering Project%
}

\ProjectTitle{%
    EduOS: A Linux-Based Academic Computing Platform%
}

\ProjectSubTitle{%
    Mid-Semester Prototype Report%
}

\author{%
    \texttt{1024170276} Dhawal Singh Jamwal \\
    \texttt{1024170275} Aaditya Kumar Yadav \\
    \texttt{1024170283} Harshit%
}

\TitlePageSubText{%
    Bachelor of Engineering - Computer Science and Engineering%
}

\AdvisorName{%
    Dr. Nisha Thakur%
}

\TitlePageFooterText{{%
    \large Mid-Semester Evaluation \\[0.2cm]
    \bfseries Computer Science and Engineering Department%
    } \\[0.15cm]
    Thapar Institute of Engineering and Technology, Patiala%
}


% ==================================================
% DOCUMENT
% ==================================================

\begin{document}

\maketitle

\tableofcontents


% ==================================================
% CHAPTER 1
% ==================================================

\chapter{Introduction}

\section{Background and Context}

University computing laboratories are required to support a wide
variety of academic activities including programming, operating
systems, networking, databases, and software development.

Although laboratory computers may contain similar hardware, their
software environments can become inconsistent because of package
changes, configuration differences, dependency problems, or changes
made by previous users.

This creates difficulties for both students and laboratory
administrators. Students may encounter systems that behave
differently from one another, while administrators must repeatedly
install, configure, repair, and maintain software across multiple
machines.

EduOS is proposed as an academic computing platform in which the
operating environment is controlled and reproducible. Instead of
treating every laboratory machine as an independently configured
computer, EduOS aims to provide a common managed academic
environment.

The current mid-semester prototype focuses on the client-side
foundation of EduOS. A custom Arch Linux image has been created
together with a multi-user student environment and a personalized
student application.

\subsection{Problem Statement}

Traditional university laboratories commonly face problems such as:

\begin{itemize}
    \item Differences in installed software between machines.
    \item Missing or incompatible development tools.
    \item Configuration changes caused by previous users.
    \item Repeated manual setup and maintenance.
    \item Lack of a common personalized academic environment.
    \item Difficulty in associating a student's authenticated
          operating-system session with academic tasks.
\end{itemize}

EduOS attempts to address these problems by providing a controlled
Linux-based academic environment that can be reproduced and extended
across laboratory machines.

\section{Project Objectives}

The major objectives of EduOS are:

\begin{enumerate}

    \item Develop a bootable customized Linux environment suitable
    for academic laboratory use.

    \item Provide separate Linux identities and home directories for
    individual students.

    \item Automatically configure required student accounts during
    system startup.

    \item Provide a personalized student interface based on the
    authenticated Linux identity.

    \item Maintain a reproducible system configuration through a
    version-controlled ArchISO profile.

    \item Design the platform so that centralized management,
    laboratory profiles, task distribution, submissions, and
    automated evaluation can be integrated in later stages.

\end{enumerate}


% ==================================================
% CHAPTER 2
% ==================================================

\chapter{Methodology and Design}

\section{System Architecture}

EduOS is designed as a modular academic computing platform.

The complete proposed architecture can be divided into four major
layers:

\begin{itemize}

    \item \textbf{EduOS Client Layer:}
    The Linux operating environment used by students on laboratory
    machines.

    \item \textbf{Management Layer:}
    Responsible for users, machines, configuration profiles, and
    centralized administration.

    \item \textbf{EduLab Layer:}
    Responsible for academic tasks, submissions, evaluation, and
    results.

    \item \textbf{Infrastructure Layer:}
    Provides supporting services such as databases, virtualization,
    containers, and deployment infrastructure.

\end{itemize}

At the current prototype stage, most implemented functionality is
within the EduOS Client Layer.

\subsection{High-Level Design}

The current prototype follows the workflow:

\begin{center}
EduOS Boot \\[2mm]
$\downarrow$ \\[2mm]
Automatic System Setup \\[2mm]
$\downarrow$ \\[2mm]
SDDM Authentication \\[2mm]
$\downarrow$ \\[2mm]
XFCE X11 Student Session \\[2mm]
$\downarrow$ \\[2mm]
Automatic EduOS Application Startup \\[2mm]
$\downarrow$ \\[2mm]
Student-Specific Academic Environment
\end{center}

The student's Linux account acts as the primary identity.

After successful authentication, the EduOS application obtains the
currently logged-in username directly from the operating system.
Therefore, a second application-level login is not required.

The username is currently mapped to academic information stored in a
local structured JSON file. In later stages, this local data source
can be replaced by a centralized API and database.

\subsection{Technical Stack}

The technologies currently used in the prototype are shown in
Table~\ref{tab:technical-stack}.

\begin{table}[H]
\centering
\renewcommand{\arraystretch}{1.25}

\begin{tabular}{|p{0.28\textwidth}|p{0.62\textwidth}|}
\hline
\textbf{Technology} & \textbf{Purpose} \\
\hline

Arch Linux &
Base operating system used by the current prototype. \\
\hline

ArchISO &
Creation of the customized bootable EduOS image. \\
\hline

XFCE &
Desktop environment used for graphical student and administrator
sessions. \\
\hline

SDDM &
Graphical display manager and Linux authentication interface. \\
\hline

Xorg / X11 &
Graphical display/session environment. \\
\hline

Python &
Implementation language for the student-side EduOS application. \\
\hline

Tkinter &
Graphical toolkit used for the current student interface. \\
\hline

JSON &
Temporary structured storage for student and task information. \\
\hline

systemd &
Boot-time execution and service management. \\
\hline

QEMU/KVM &
Virtualization environment used to test generated ISO images. \\
\hline

Git &
Version control for source code and ArchISO configuration. \\
\hline

FastAPI &
Planned technology for the centralized management API. \\
\hline

PostgreSQL &
Planned persistent data storage for future backend services. \\
\hline

Podman &
Planned isolation mechanism for automated evaluation workloads. \\
\hline

\end{tabular}

\caption{EduOS Technical Stack}
\label{tab:technical-stack}
\end{table}


% ==================================================
% IMPLEMENTATION DETAILS
% ==================================================

\section{Implementation Details}


% --------------------------------------------------
% MODULE 1
% --------------------------------------------------

\subsection{Module 1: Arch Linux ISO Setup}

The first implementation module focuses on creating the base EduOS
operating-system environment using Arch Linux and ArchISO.

The purpose of this module is to produce a bootable customized Linux
image containing the graphical environment, runtime dependencies,
system services, and configuration required by EduOS.

The ArchISO profile is maintained under:

\begin{verbatim}
code/archiso/profile
\end{verbatim}

The current image contains:

\begin{itemize}

    \item Arch Linux base system.
    \item XFCE desktop environment.
    \item SDDM display manager.
    \item Xorg graphical environment.
    \item Mesa graphics support.
    \item Python runtime.
    \item Tkinter.
    \item systemd.
    \item Networking and basic system administration utilities.

\end{itemize}

The ISO is generated using the \texttt{mkarchiso} utility.

The general module workflow is:

\begin{center}
ArchISO Profile \\[2mm]
$\downarrow$ \\[2mm]
Package Selection \\[2mm]
$\downarrow$ \\[2mm]
Configuration Integration \\[2mm]
$\downarrow$ \\[2mm]
ISO Generation \\[2mm]
$\downarrow$ \\[2mm]
QEMU/KVM Testing
\end{center}

QEMU/KVM allows the complete EduOS environment to be tested in an
isolated virtual machine without modifying the host operating
system.

During implementation, an important issue occurred in the graphical
session startup. SDDM authentication succeeded, but the XFCE session
terminated because the Wayland XFCE session was being selected.

The problem was resolved by explicitly configuring SDDM to use the
X11 graphical environment.

The relevant configuration is stored at:

\begin{verbatim}
/etc/sddm.conf.d/10-eduos-x11.conf
\end{verbatim}

After this configuration was introduced, SDDM could authenticate
users and successfully launch the XFCE X11 session.

This module therefore provides the bootable and reproducible
operating-system foundation for EduOS.


% --------------------------------------------------
% MODULE 2
% --------------------------------------------------

\subsection{Module 2: User Base and Student Environment Setup}

The second module focuses on establishing a multi-user academic
environment on top of the EduOS operating system.

Instead of representing students only inside an application,
students are implemented as genuine Linux user accounts.

The current prototype contains five test student accounts:

\begin{itemize}
    \item \texttt{student1000}
    \item \texttt{student1001}
    \item \texttt{student1002}
    \item \texttt{student1003}
    \item \texttt{student1004}
\end{itemize}

A separate \texttt{eduos} account is retained for system
administration, maintenance, and debugging.

Each student receives:

\begin{itemize}
    \item A separate Linux username.
    \item A unique Linux user identity.
    \item An independent home directory.
    \item A separate login session.
    \item A configured password.
    \item An associated roll number.
    \item Student-specific academic information.
\end{itemize}

Student account creation is automated using:

\begin{verbatim}
/usr/local/bin/eduos-setup
\end{verbatim}

The setup script is executed by a systemd service during system
startup.

The service runs before the graphical display manager, ensuring that
the required student accounts exist before SDDM presents the login
screen.

The provisioning process is designed to be idempotent. Re-running
the setup process maintains the required user configuration without
creating duplicate student accounts.

The student interface is implemented in Python using Tkinter:

\begin{verbatim}
/usr/local/bin/eduos-student-app
\end{verbatim}

After login, the application retrieves the currently authenticated
Linux username.

Student academic information is currently read from:

\begin{verbatim}
/usr/share/eduos/student-data.json
\end{verbatim}

The interface displays information including:

\begin{itemize}
    \item Student name.
    \item Roll number.
    \item Linux username.
    \item Assigned academic tasks.
    \item Task status.
    \item Due information.
\end{itemize}

Different test students are assigned different tasks to demonstrate
that the environment is personalized according to the authenticated
identity.

The application launches automatically after student login and opens
in fullscreen mode.

The administrative \texttt{eduos} account is excluded from this
student workflow and instead receives access to the normal XFCE
desktop.


% ==================================================
% SYSTEM DIAGRAMS
% ==================================================

\section{System Diagrams}


% --------------------------------------------------
% USE CASE
% --------------------------------------------------

\subsection{Use Case Diagram}

The use case diagram represents the major interactions between
students, administrators, and the EduOS system.

\diagramfigure
{images/use-case-diagram.png}
{EduOS Use Case Diagram}
{fig:use-case}


% --------------------------------------------------
% USE CASE SPECIFICATIONS
% --------------------------------------------------

\subsection{Use Case Specifications}

The use case specifications describe important system interactions
in greater detail by defining actors, goals, preconditions, normal
flows, alternative flows, and postconditions.

\diagramfigure
{images/use-case-specifications.png}
{EduOS Use Case Specifications}
{fig:use-case-specifications}


% --------------------------------------------------
% CLASS DIAGRAM
% --------------------------------------------------

\subsection{Class Diagram}

The class diagram represents the major conceptual entities of EduOS
and the relationships between users, students, administrators,
tasks, submissions, results, machines, and supporting system
components.

\diagramfigure
{images/class-diagram.png}
{EduOS Class Diagram}
{fig:class-diagram}


% --------------------------------------------------
% ACTIVITY / SWIMLANE
% --------------------------------------------------

\subsection{Activity Diagram}

The activity diagram represents the workflow of the EduOS system.
Swimlanes are used to distinguish the responsibilities of the
student and the major EduOS components.

\diagramfigure
{images/activity-diagram.png}
{EduOS Activity / Swimlane Diagram}
{fig:activity}


% --------------------------------------------------
% DATA FLOW DIAGRAMS
% --------------------------------------------------

\subsection{Data Flow Diagrams}

Data Flow Diagrams are used to represent how information moves
between users, EduOS processes, and data stores.


\subsubsection{Level 0 Data Flow Diagram}

The Level 0 DFD provides a context-level representation of EduOS and
its interactions with external entities.

\diagramfigure
{images/dfd-level-0.png}
{EduOS Level 0 Data Flow Diagram}
{fig:dfd0}


\subsubsection{Level 1 Data Flow Diagram}

The Level 1 DFD decomposes EduOS into its major functional
processes, including authentication, machine management, task
management, evaluation, and results management.

\diagramfigure
{images/dfd-level-1.png}
{EduOS Level 1 Data Flow Diagram}
{fig:dfd1}


\subsubsection{Level 2 Data Flow Diagram}

The Level 2 DFD provides a more detailed decomposition of the task
and submission workflow.

\diagramfigure
{images/dfd-level-2.png}
{EduOS Level 2 Data Flow Diagram}
{fig:dfd2}


% --------------------------------------------------
% SEQUENCE
% --------------------------------------------------

\subsection{Sequence Diagram}

The sequence diagram represents the chronological exchange of
messages between the student and major EduOS components.

\diagramfigure
{images/sequence-diagram.png}
{EduOS Sequence Diagram}
{fig:sequence}


% --------------------------------------------------
% COLLABORATION
% --------------------------------------------------

\subsection{Collaboration Diagram}

The collaboration diagram represents the structural interaction
between major EduOS components and emphasizes their communication
relationships.

\diagramfigure
{images/collaboration-diagram.png}
{EduOS Collaboration Diagram}
{fig:collaboration}


% --------------------------------------------------
% GANTT
% --------------------------------------------------

\subsection{Gantt Chart}

The Gantt chart represents the development schedule and planned
sequence of major EduOS activities.

\diagramfigure
{images/gantt-chart.png}
{EduOS Project Development Timeline}
{fig:gantt}


% ==================================================
% CHAPTER 3
% ==================================================

\chapter{Results and Evaluation}

\section{Testing Methodology}

The current EduOS prototype was evaluated primarily through
functional and integration testing.

Since the project is currently at the proof-of-concept stage, the
main objective of testing was to verify whether the operating-system
image, authentication process, student sessions, automated setup,
and student application operate correctly as a complete workflow.

The following testing approaches were used:

\begin{itemize}

    \item \textbf{Component Testing:}
    Individual scripts and configuration files were checked before
    being included in the ISO. This included the account provisioning
    script, JSON student data, application startup configuration, and
    Python student application.

    \item \textbf{Integration Testing:}
    The interaction between systemd, SDDM, XFCE, Linux users, and the
    EduOS student application was tested after generating the complete
    system image.

    \item \textbf{Virtual Machine Testing:}
    Generated ISO images were booted using QEMU/KVM to test the system
    from a clean startup environment.

    \item \textbf{Multi-User Testing:}
    Different student accounts were authenticated to confirm that
    users receive separate Linux identities, home directories, and
    student-specific academic information.

    \item \textbf{Administrative Testing:}
    The \texttt{eduos} account was tested separately to verify that
    administrative access to the normal XFCE environment remains
    available.

\end{itemize}


\section{Performance Metrics}

At the current prototype stage, evaluation focuses primarily on
functional correctness rather than large-scale performance
benchmarking.

Table~\ref{tab:test-results} summarizes the major prototype
validation results.

\begin{table}[H]
\centering
\renewcommand{\arraystretch}{1.25}

\begin{tabular}{|p{0.46\textwidth}|p{0.18\textwidth}|p{0.22\textwidth}|}
\hline
\textbf{Test} & \textbf{Expected} & \textbf{Result} \\
\hline

Custom EduOS ISO boots in QEMU/KVM &
Successful boot &
Passed \\
\hline

SDDM graphical login is displayed &
Login available &
Passed \\
\hline

Student accounts authenticate &
Valid login &
Passed \\
\hline

Five separate student Linux accounts exist &
5 users &
Passed \\
\hline

Separate student home directories are created &
Independent homes &
Passed \\
\hline

XFCE X11 session starts after authentication &
Desktop session &
Passed \\
\hline

EduOS student application starts automatically &
Automatic startup &
Passed \\
\hline

Application detects authenticated Linux user &
Correct identity &
Passed \\
\hline

Different users receive different task information &
Personalization &
Passed \\
\hline

Student application opens in fullscreen mode &
Fullscreen interface &
Passed \\
\hline

Administrative account retains normal XFCE desktop &
Admin access &
Passed \\
\hline

\end{tabular}

\caption{Prototype Functional Evaluation Results}
\label{tab:test-results}
\end{table}


\section{Analysis}

The prototype successfully demonstrates the primary client-side
concept of EduOS.

A customized Linux image can be generated and booted inside a virtual
machine. Student accounts are created at the operating-system level,
allowing standard Linux authentication and session isolation to be
used instead of implementing a separate simulated authentication
system.

The system also demonstrates identity-based personalization. Once a
student has authenticated, the EduOS application can determine the
current Linux user and retrieve the corresponding academic data.

This demonstrates a complete prototype flow:

\begin{center}
Boot \\[2mm]
$\downarrow$ \\[2mm]
Authenticate \\[2mm]
$\downarrow$ \\[2mm]
Create Student Session \\[2mm]
$\downarrow$ \\[2mm]
Identify Student \\[2mm]
$\downarrow$ \\[2mm]
Load Personalized Academic Data
\end{center}

Several implementation challenges were encountered during
development.

The most significant graphical-session issue occurred when the
display manager attempted to launch XFCE using Wayland. The problem
was diagnosed through session behavior and configuration inspection,
and the system was configured to use X11 instead.

Another important challenge was ensuring that the student accounts
were available before graphical authentication. This was addressed
by integrating account provisioning with systemd and controlling the
startup order relative to SDDM.

The current results therefore demonstrate that the underlying EduOS
client architecture is technically feasible.

However, the current academic data is still stored locally and the
prototype does not yet provide the centralized management,
submission, or automated evaluation components planned for the full
system.


% ==================================================
% CHAPTER 4
% ==================================================

\chapter{Conclusion}

\section{Summary of Work}

The mid-semester EduOS prototype establishes the foundation of a
Linux-based academic computing environment.

Two major implementation modules were completed during the current
stage.

The first module produced a customized bootable Arch Linux image
using ArchISO. The image integrates XFCE, SDDM, Xorg, Python,
Tkinter, systemd, and the supporting packages required by the
current prototype.

The second module implemented a multi-user student environment.
Multiple real Linux student accounts are automatically provisioned,
each with an independent identity, home directory, login session,
and associated academic information.

A student-side Python/Tkinter application was also integrated into
the operating system. It automatically identifies the currently
authenticated Linux user and displays student-specific information
and assigned tasks.

Together, these modules demonstrate the central client-side concept
of EduOS: a student can boot into a standardized academic operating
environment, authenticate using a Linux account, and automatically
receive a personalized workspace.


\section{Key Findings}

The major findings from the current implementation are:

\begin{itemize}

    \item \textbf{Linux-native authentication can be reused for
    academic identity.}
    The student application does not require a second login because
    it can use the identity of the authenticated Linux session.

    \item \textbf{Student environments can be provisioned
    automatically.}
    systemd-based initialization allows user accounts and required
    configuration to be established during system startup.

    \item \textbf{A customized academic Linux image can be reproduced
    using ArchISO.}
    Package selection and system configuration are maintained within
    the project profile rather than being manually applied after
    installation.

    \item \textbf{Student-specific interfaces can coexist with an
    administrative environment.}
    Student accounts automatically enter the EduOS interface, while
    the administrative account retains access to normal XFCE.

    \item \textbf{Virtual-machine testing is effective during
    operating-system development.}
    QEMU/KVM allowed repeated ISO builds and boot testing without
    modifying the host machine.

\end{itemize}


\section{Future Work and Recommendations}

The next development stage will extend the current client-side
prototype into the broader EduOS architecture.

Planned work includes:

\begin{enumerate}

    \item Develop a centralized management API using FastAPI.

    \item Replace the local JSON student data with persistent
    database-backed information.

    \item Implement centralized student and machine management.

    \item Introduce configurable laboratory profiles for different
    courses and practical sessions.

    \item Implement task distribution and submission handling.

    \item Develop the first automated programming evaluator.

    \item Introduce isolated execution using containers.

    \item Store and present evaluation results and feedback.

    \item Connect EduOS clients to the management server so that
    laboratory machines can receive centrally controlled
    configurations.

\end{enumerate}


\section{Final Remarks}

The current prototype demonstrates that EduOS can provide more than
a conventional customized Linux desktop.

By combining a reproducible operating-system image, native Linux
user management, automated provisioning, and a student-specific
interface, the prototype establishes a foundation for a managed
academic computing platform.

The next stages of development will focus on moving from the current
local proof of concept toward centralized management and automated
academic assessment.

\end{document}
